\section{Introduction}
% ==================================================================================================================

Rocket nozzle geometry is one of the dominant factors controlling the conversion of chamber thermal energy into directed exhaust kinetic energy. For a fixed chamber pressure, propellant mass flow rate and propellant combination, the nozzle determines the exhaust Mach number, exit pressure, exit velocity, thrust coefficient and degree of axial momentum alignment. An efficient nozzle must therefore expand the combustion gases as much as possible while avoiding excessive length, excessive divergence losses, flow separation, shock formation or manufacturing complexity.

Among the most common rocket nozzle configurations, the bell nozzle offers a strong compromise between performance and compactness. A purely conical nozzle is simple to manufacture and analyze, but its constant wall angle produces larger divergence losses for a given length. A bell nozzle replaces most of the conical diverging section with a curved contour, allowing strong initial expansion immediately downstream of the throat and a smaller exit angle at the nozzle exit. This improves axial momentum alignment while maintaining a shorter nozzle than an equivalent low-angle conical design.

The simulator described in this work was created to generate and analyze single bell nozzle geometries for hybrid rocket engines, with emphasis on a paraffin wax and nitrous oxide propulsion system. The original work mixed implementation details, code excerpts, theoretical derivations, geometry construction and results in the same sections. This revised manuscript separates the material into four layers: physical theory, geometric formulation, numerical implementation logic, and results/optimization interpretation.

The simulator is intended as a preliminary engineering tool rather than a final qualification model. Its value lies in the explicit connection between geometry, thermochemical properties, quasi-one-dimensional flow, viscous-loss closures, visualization, CAD export and performance trends. The current implementation integrates these models with a genetic algorithm so that the expansion ratio and divergent-contour parameters can be optimized automatically at a specified operating point.
